\FigureLayoutDeclare{img/2010/beijing_liberal/q5_condition_front_fixed.png}{4.4cm}{11bp 0bp 13bp 5bp}
\FigureLayoutDeclare{img/2010/beijing_liberal/q5_condition_side_fixed.png}{2.5cm}{10bp 0bp 11bp 5bp}
\FigureLayoutDeclare{img/2010/beijing_liberal/q7_fig1.png}{3.5cm}{3bp 28bp 31bp 0bp}

\chapter{2010年北京卷（文）}

\section{单选题}

\begin{problem}
集合 \(P = \{x \in \Z \mid 0 \leqslant x < 3\}\)，\(M = \{x \in \Z \mid x^2 \leqslant 9\}\)，则 \(P \cap M =\)（\quad）

\choices
    {\( \{1, 2\} \)}
    {\(\{0, 1, 2\}\)}
    {\(\{1, 2, 3\}\)}
    {\(\{0, 1, 2, 3\}\)}
\end{problem}

\begin{answer}
B
\end{answer}

\begin{solution}
由题意，\(0\leqslant x<3\) 且 \(x\in\Z\)，所以 \(P=\{0,1,2\}\).又因为 \(x^2\leqslant9\)，所以 \(-3\leqslant x\leqslant3\)，从而 \(M=\{-3,-2,-1,0,1,2,3\}\).因此
\[
P\cap M=\{0,1,2\}
\]
故选 B.
\end{solution}

\begin{problem}
在复平面内，复数 \(6 + 5\i\)，\(-2 + 3\i\) 对应的点分别为 \(A\)，\(B\). 若 \(C\) 为线段 \(AB\) 的中点，则点 \(C\) 对应的复数是（\quad）

\choices
    {\(4 + 8\i\)}
    {\(8 + 2\i\)}
    {\(2 + 4\i\)}
    {\(4 + \i\)}
\end{problem}

\begin{answer}
C
\end{answer}

\begin{solution}
复数 \(6+5\i\)，\(-2+3\i\) 对应的点的坐标分别为 \((6,5)\)，\((-2,3)\).线段中点 \(C\) 的坐标为
\[
\left(\frac{6+(-2)}{2},\frac{5+3}{2}\right)=(2,4)
\]
所以点 \(C\) 对应的复数为 \(2+4\i\)，故选 C.
\end{solution}

\begin{problem}
从 \(\{1,2,3,4,5\}\) 中随机选取一个数为 \(a\)，从 \(\{1,2,3\}\) 中随机选取一个数为 \(b\)，则 \(b > a\) 的概率是（\quad）

\choices
    {\( \frac{4}{5} \)}
    {\( \frac{3}{5} \)}
    {\( \frac{2}{5} \)}
    {\( \frac{1}{5} \)}
\end{problem}

\begin{answer}
D
\end{answer}

\begin{solution}
从两组数中各取一个数的基本结果共有 \(5\times3=15\) 个，且等可能.满足 \(b>a\) 的结果为
\((a,b)=(1,2)\)，\((1,3)\)，\((2,3)\)
共 \(3\) 个.因此
\[
P(b>a)=\frac{3}{15}=\frac{1}{5}
\]
故选 D.
\end{solution}

\begin{problem}
若 \(\bs{a}\)，\(\bs{b}\) 是非零向量，且 \(\bs{a} \perp \bs{b}\)，\(|\bs{a}| \neq |\bs{b}|\)，则函数 \(f(x) = (x\bs{a} + \bs{b}) \cdot (x\bs{b} - \bs{a})\) 是（\quad）

\choices
    {一次函数且是奇函数}
    {一次函数但不是奇函数}
    {二次函数且是偶函数}
    {二次函数但不是偶函数}
\end{problem}

\begin{answer}
A
\end{answer}

\begin{solution}
由 \(\bs{a}\perp\bs{b}\)，得 \(\bs{a}\cdot\bs{b}=0\).展开数量积：
\[
f(x)=(x\bs{a}+\bs{b})\cdot(x\bs{b}-\bs{a})
=x^2\bs{a}\cdot\bs{b}-x|\bs{a}|^2+x|\bs{b}|^2-\bs{b}\cdot\bs{a}
=(|\bs{b}|^2-|\bs{a}|^2)x
\]
因为 \(|\bs{a}|\ne|\bs{b}|\)，一次项系数非零，所以 \(f(x)\) 是一次函数；又 \(f(-x)=-f(x)\)，所以它是奇函数，故选 A.
\end{solution}

\begin{problem}
一个长方体去掉一个小长方体，所得几何体的正（主）视图与侧（左）视图分别如图所示，则该几何体的俯视图为（\quad）

\examfiguregroup[float=inline]{img/2010/beijing_liberal/q5_condition_front_fixed.png}{%
    \bitmapinclude[width=0.26\linewidth]{img/2010/beijing_liberal/q5_condition_front_fixed.png}\hspace{1.2em}%
    \bitmapinclude[width=0.15\linewidth]{img/2010/beijing_liberal/q5_condition_side_fixed.png}%
}

\choices
    {\choicebitmap[width=0.15\paperwidth]{img/2010/beijing_liberal/q5_optA_fixed.png}}
    {\choicebitmap[width=0.15\paperwidth]{img/2010/beijing_liberal/q5_optB_fixed.png}}
    {\choicebitmap[width=0.15\paperwidth]{img/2010/beijing_liberal/q5_optC_fixed.png}}
    {\choicebitmap[width=0.15\paperwidth]{img/2010/beijing_liberal/q5_optD_fixed.png}}
\end{problem}

\begin{answer}
C
\end{answer}

\begin{solution}
把原长方体的长度、宽度和高度分别看作正视图的横向、侧视图的横向和两个视图共同的竖向.正视图中缺口位于左上方，说明被去掉的小长方体在长度方向靠左、在高度方向靠上；侧视图中缺口位于右上方，说明它在宽度方向所对应的位置确定在俯视图的下方.因此，从上方观察时，缺口应位于外框的左下方.四个选项中只有 C 的缺口位置同时满足这两个投影条件，故选 C.
\end{solution}

\begin{problem}
给定函数：

\begin{circlelist}
\item \( y = x^{2} \)
\item \( y = \log_{\frac{1}{2}}(x + 1) \)
\item \( y = |x - 1| \)
\item \( y = 2^{x + 1} \)
\end{circlelist}

其中在区间 \( (0, 1) \) 上单调递减的函数序号是（\quad）

\choices
    {\circled{1}\circled{2}}
    {\circled{2}\circled{3}}
    {\circled{3}\circled{4}}
    {\circled{1}\circled{4}}
\end{problem}

\begin{answer}
B
\end{answer}

\begin{solution}
在 \((0,1)\) 上，\(y=x^2\) 随 \(x\) 增大而增大，不是减函数；由于底数 \(\frac{1}{2}\) 小于 \(1\)，且 \(x+1\) 随 \(x\) 增大而增大，所以 \(y=\log_{\frac{1}{2}}(x+1)\) 单调递减；\(y=|x-1|=1-x\) 在 \((0,1)\) 上单调递减；\(y=2^{x+1}\) 单调递增.因此单调递减的函数序号为 \(\circled{2}\circled{3}\)，故选 B.
\end{solution}

\begin{problem}
某班设计了一个八边形的班徽（如图），它由腰长为 1，顶角为 \(\alpha\) 的四个等腰三角形，及其底边构成的正方形所组成，该八边形的面积为（\quad）

\bitmapfigure[max width=0.35\linewidth]{img/2010/beijing_liberal/q7_fig1.png}

\choices
    {\(2\sin \alpha - 2\cos \alpha + 2\)}
    {\(\sin \alpha -\sqrt{3}\cos \alpha +3\)}
    {\(3\sin \alpha -\sqrt{3}\cos \alpha +1\)}
    {\(2\sin \alpha -\cos \alpha +1\)}
\end{problem}

\begin{answer}
A
\end{answer}

\begin{solution}
每个等腰三角形的两腰均为 \(1\)，顶角为 \(\alpha\)，所以一个三角形的面积为
\[
\frac{1}{2}\times1\times1\times\sin\alpha=\frac{1}{2}\sin\alpha
\]
四个三角形的面积和为 \(2\sin\alpha\).由余弦定理，正方形边长 \(s\) 满足
\[
s^2=1^2+1^2-2\cos\alpha=2-2\cos\alpha
\]
所以正方形面积为 \(2-2\cos\alpha\).八边形面积为
\[
2\sin\alpha+2-2\cos\alpha=2\sin\alpha-2\cos\alpha+2
\]
故选 A.
\end{solution}

\begin{problem}
如图，正方体 \(ABCD - A_{1}B_{1}C_{1}D_{1}\) 的棱长为 2 ，动点 \(E\)，\(F\) 在棱 \(A_{1}B_{1}\) 上. 点 \(Q\) 是棱 \(CD\) 的中点，动点 \(P\) 在棱 \(AD\) 上，若 \(EF = 1\)，\(DP = x\)，\(A_{1}E = y\)（\(x\)，\(y\) 大于零），则三棱锥 \(P - EFQ\) 的体积（\quad）

\bitmapfigure[max width=0.34\linewidth]{img/2010/beijing_liberal/q8_fig1.png}

\choices
    {与 \(x, y\) 都有关}
    {与 \(x\)，\(y\) 都无关}
    {与 \(x\) 有关，与 \(y\) 无关}
    {与 \(y\) 有关，与 \(x\) 无关}
\end{problem}

\begin{answer}
C
\end{answer}

\begin{solution}
建立直角坐标系，取 \(A(0,0,0)\)，\(B(2,0,0)\)，\(C(2,2,0)\)，\(D(0,2,0)\)，\(A_1(0,0,2)\)，\(B_1(2,0,2)\).设 \(P=(0,2-x,0)\)，\(E=(y,0,2)\)，\(F=(y\pm1,0,2)\)，其中符号取决于 \(E\)，\(F\) 在棱上的先后次序，且 \(EF=1\).又 \(Q\) 是 \(CD\) 的中点，所以 \(Q=(1,2,0)\).

平面 \(EFQ\) 的方程为 \(Y+Z=2\)，与 \(y\) 无关.三角形 \(EFQ\) 的底边 \(EF=1\)，高为
\[
\sqrt{2^2+(-2)^2}=2\sqrt{2}
\]
故 \(S_{\triangle EFQ}=\sqrt{2}\).点 \(P\) 到平面 \(EFQ\) 的距离为
\[
\frac{|(2-x)+0-2|}{\sqrt{1^2+1^2}}=\frac{x}{\sqrt{2}}
\]
因此
\[
V_{P-EFQ}=\frac{1}{3}\times\sqrt{2}\times\frac{x}{\sqrt{2}}=\frac{x}{3}
\]
只与 \(x\) 有关，与 \(y\) 无关，故选 C.
\end{solution}

\section{填空题}

\begin{problem}
已知函数 \( y = \begin{cases}
\log_{2} x, & x \geqslant 2\\
2 - x, & x < 2
\end{cases} \)，如图表示的是给定 \(x\) 的值，求其对应的函数值 \(y\) 的程序框图. \circled{1} 应填写\fillinblank{}；\circled{2} 应填写\fillinblank{}.

\texfigure[max width=0.35\linewidth]{img_repaint/2010/beijing_liberal/q9_fig1.tex}
\end{problem}

\begin{answer}
\(x < 2\)；\(y = \log_{2}x\)
\end{answer}

\begin{solution}
流程图在判断框中“是”的分支直接执行 \(y=2-x\)，所以该分支对应分段函数的条件 \(x<2\).判断框中“否”的分支对应 \(x\geqslant2\)，此时应执行另一段 \(y=\log_{2}x\).因此 \(\circled{1}\) 填 \(x<2\)，\(\circled{2}\) 填 \(y=\log_{2}x\).
\end{solution}

\begin{problem}
在 \(\triangle ABC\) 中，若 \(b = 1\)，\(c = \sqrt{3}\)，\(\angle C = \frac{2\pi}{3}\)，则 \(a = \)\fillinblank{}.
\end{problem}

\begin{answer}
\(1\)
\end{answer}

\begin{solution}
由正弦定理，
\[
\frac{b}{\sin B}=\frac{c}{\sin C}
\]
所以
\[
\sin B=\frac{b\sin C}{c}=\frac{1\times\sin\frac{2\pi}{3}}{\sqrt{3}}=\frac{1}{2}
\]
故 \(B=\frac{\pi}{6}\) 或 \(B=\frac{5\pi}{6}\).但 \(B+C<\pi\)，而 \(C=\frac{2\pi}{3}\)，所以 \(B<\frac{\pi}{3}\)，只能取 \(B=\frac{\pi}{6}\).于是 \(A=\frac{\pi}{6}\).再次由正弦定理，
\[
a=\frac{c\sin A}{\sin C}=\frac{\sqrt{3}\times\frac{1}{2}}{\frac{\sqrt{3}}{2}}=1
\]
故填 \(1\).
\end{solution}

\begin{problem}
若点 \(P(m,3)\) 到直线 \(4x - 3y + 1 = 0\) 的距离为4，且点 \(P\) 在不等式 \(2x + y < 3\) 表示的平面区域内，则 \(m =\fillinblank{}\).
\end{problem}

\begin{answer}
\(-3\)
\end{answer}

\begin{solution}
由点到直线的距离公式，
\[
\frac{|4m-3\times3+1|}{\sqrt{4^2+(-3)^2}}=4
\]
即 \(|4m-8|=20\).因此 \(m=7\) 或 \(m=-3\).又点 \(P(m,3)\) 在 \(2x+y<3\) 表示的区域内，故 \(2m+3<3\)，即 \(m<0\)，所以只能取 \(m=-3\).
\end{solution}

\begin{problem}
从某小学随机抽取 100 名同学，将他们的身高（单位：厘米）数据绘制成频率分布直方图（如图）.
由图中数据可知 \(a =\)\fillinblank{}. 若要从身高在 \([120, 130)\)，\([130, 140)\)，\([140, 150]\) 三组内的学生中，用分层抽样的方法选取 18 人参加一项活动，则从身高在 \([140, 150]\) 内的学生中选取的人数应为\fillinblank{}.
\texfigure[max width=0.55\linewidth]{img_repaint/2010/beijing_liberal/q12_fig1.tex}
\end{problem}

\begin{answer}
\(0.030\)，\(3\)
\end{answer}

\begin{solution}
频率分布直方图中各小矩形面积之和为 \(1\)，所以
\[
10(0.005+0.035+a+0.020+0.010)=1
\]
解得 \(a=0.030\).

身高在 \([120,130)\)，\([130,140)\)，\([140,150]\) 三组内的总频率为
\[
10(0.030+0.020+0.010)=0.6
\]
故这三组共有 \(100\times0.6=60\) 人.其中 \([140,150]\) 组的频率为 \(10\times0.010=0.1\)，共有 \(10\) 人.按分层抽样比例，应选取
\[
18\times\frac{10}{60}=3
\]
人.因此两空分别为 \(0.030\) 和 \(3\).
\end{solution}

\begin{problem}
已知双曲线 \(\frac{x^2}{a^2} - \frac{y^2}{b^2} = 1\)（\(a > 0\)，\(b > 0\)）的离心率为 2，焦点与椭圆 \(\frac{x^2}{25} + \frac{y^2}{9} = 1\) 的焦点相同，那么双曲线的焦点坐标为\fillinblank{}；渐近线方程为\fillinblank{}.
\end{problem}

\begin{answer}
\((4,0)\)，\((-4,0)\)；\(y = \pm\sqrt{3}x\)
\end{answer}

\begin{solution}
椭圆 \(\dfrac{x^2}{25}+\dfrac{y^2}{9}=1\) 的焦半距满足
\[
c^2=25-9=16
\]
所以焦点为 \((4,0)\)，\((-4,0)\).双曲线与该椭圆的焦点相同，故其焦半距也是 \(c=4\).又双曲线离心率为 \(2\)，所以
\(\frac{c}{a}=2\)，\(a=2\).
由 \(c^2=a^2+b^2\)，得 \(b^2=16-4=12\).双曲线的渐近线为
\[
y=\pm\frac{b}{a}x=\pm\sqrt{3}x
\]
故填 \((4,0)\)，\((-4,0)\) 和 \(y=\pm\sqrt{3}x\).
\end{solution}

\begin{problem}
如图放置的边长为 1 的正方形 \(PABC\) 沿 \(x\) 轴滚动. 设顶点 \(P(x, y)\) 的纵坐标与横坐标的函数关系是 \(y = f(x)\)，则 \(f(x)\) 的最小正周期为\fillinblank{}；\(y = f(x)\) 在其两个相邻零点间的图象与 \(x\) 轴所围区域的面积为\fillinblank{}.

说明：“正方形 \(PABC\) 沿 \(x\) 轴滚动”包括沿 \(x\) 轴正方向和沿 \(x\) 轴负方向滚动. 沿 \(x\) 轴正方向滚动指的是先以顶点 \(A\) 为中心顺时针旋转，当顶点 \(B\) 落在 \(x\) 轴上时，再以顶点 \(B\) 为中心顺时针旋转，如此继续. 类似地，正方形 \(PABC\) 可以沿 \(x\) 轴负方向滚动.

\bitmapfigure[max width=0.32\linewidth]{img/2010/beijing_liberal/q14_fig1.png}
\end{problem}

\begin{answer}
\(4\)；\(\pi + 1\)
\end{answer}

\begin{solution}
正方形每次绕一个落在 \(x\) 轴上的顶点旋转 \(90^{\circ}\)，连续四次后恢复原来的姿态，并沿 \(x\) 轴平移 \(4\) 个单位.因此 \(f(x+4)=f(x)\).在一个滚动周期内，顶点 \(P\) 除周期端点外始终在 \(x\) 轴上方，所以 \(f(x)\) 的零点相邻间距均为 \(4\).任何正周期都必须把零点映射为零点，故最小正周期为 \(4\).

考察正方形向右滚动时 \(P\) 从一个零点到下一个零点的轨迹.第一次绕 \(A\) 转过 \(90^{\circ}\)，轨迹为半径 \(1\) 的四分之一圆，面积为 \(\frac{\pi}{4}\).第二次绕 \(B\) 转过 \(90^{\circ}\)，\(BP=\sqrt{2}\)，轨迹对应圆心在 \(x\) 轴上的半径 \(\sqrt{2}\) 的圆弧，横坐标从 \(1\) 到 \(3\)，其与 \(x\) 轴围成的面积为
\[
\int_{1}^{3}\sqrt{2-(x-2)^2}\,\mathrm{d}x
=\left[\frac{x-2}{2}\sqrt{2-(x-2)^2}+\arcsin\frac{x-2}{\sqrt{2}}\right]_{1}^{3}
=1+\frac{\pi}{2}
\]
第三次绕 \(C\) 转过 \(90^{\circ}\)，轨迹为半径 \(1\) 的四分之一圆，面积为 \(\frac{\pi}{4}\).所以所求面积为
\[
\frac{\pi}{4}+1+\frac{\pi}{2}+\frac{\pi}{4}=\pi+1
\]
故两空分别为 \(4\) 和 \(\pi+1\).
\end{solution}

\section{解答题}

\begin{problem}
已知函数 \(f(x)=2\cos 2x+\sin^2 x\).

\begin{enumerate}
\item 求 \( f\left(\frac{\pi}{3}\right) \) 的值；
\item 求 \( f(x) \) 的最大值和最小值.
\end{enumerate}
\end{problem}

\begin{solution}
\begin{enumerate}
\item
\[
f\left(\frac{\pi}{3}\right)=2\cos\frac{2\pi}{3}+\sin^2\frac{\pi}{3}=2\times\left(-\frac{1}{2}\right)+\left(\frac{\sqrt{3}}{2}\right)^2=-\frac{1}{4}
\]

\item 令 \(u=\cos x\)，则 \(u\in[-1,1]\).利用 \(\sin^2x=1-\cos^2x\) 和 \(\cos2x=2\cos^2x-1\)，得
\[
f(x)=2(2u^2-1)+(1-u^2)=3u^2-1
\]
因为 \(u^2\in[0,1]\)，所以 \(f(x)\in[-1,2]\).当 \(u=0\) 时取最小值 \(-1\)，当 \(u=\pm1\) 时取最大值 \(2\).
\end{enumerate}
\end{solution}

\begin{problem}
已知 \(\{a_{n}\}\) 为等差数列，且 \(a_{3} = -6\)，\(a_{6} = 0\).

\begin{enumerate}
\item 求 \( \{a_{n}\} \) 的通项公式；
\item 若等比数列 \( \{b_{n}\} \) 满足 \(b_{1} = -8\)，\(b_{2} = a_{1} + a_{2} + a_{3}\)，求 \( \{b_{n}\} \) 的前 \(n\) 项和公式.
\end{enumerate}
\end{problem}

\begin{solution}
\begin{enumerate}
\item 设等差数列 \(\{a_n\}\) 的公差为 \(d\)，则
\[
\begin{cases}
a_1+2d=-6,\\
a_1+5d=0.
\end{cases}
\]
两式相减得 \(3d=6\)，所以 \(d=2\)，进而 \(a_1=-10\).因此
\[
a_n=a_1+(n-1)d=-10+2(n-1)=2n-12
\]

\item 由上式，\(a_1=-10\)，\(a_2=-8\)，\(a_3=-6\)，所以
\[
b_2=a_1+a_2+a_3=-24
\]
设等比数列 \(\{b_n\}\) 的公比为 \(q\)，则 \(-8q=-24\)，得 \(q=3\).故其前 \(n\) 项和为
\[
S_n=b_1\frac{q^n-1}{q-1}=-8\frac{3^n-1}{2}=-4(3^n-1)
\]
\end{enumerate}
\end{solution}

\begin{problem}
如图，正方形 \(ABCD\) 和四边形 \(ACEF\) 所在的平面互相垂直，\(EF \parallel AC\)，\(AB = \sqrt{2}\)，\(CE = EF = 1\).

\begin{enumerate}
\item 求证：\( AF \parallel \) 平面 \(BDE\)；
\item 求证：\( CF \perp \) 平面 \(BDE\).
\end{enumerate}

\bitmapfigure[max width=0.44\linewidth]{img/2010/beijing_liberal/q17_fig1.png}
\end{problem}

\begin{solution}
\begin{enumerate}
\item 设 \(G\) 为 \(AC\) 与 \(BD\) 的交点.因为 \(ABCD\) 是边长为 \(\sqrt{2}\) 的正方形，所以 \(AC=2\)，从而 \(AG=CG=1\).又 \(EF\parallel AC\)，且 \(EF=AG=1\)，所以四边形 \(AGEF\) 是平行四边形，得到 \(AF\parallel GE\).由于 \(G\in BD\)，\(E\) 在平面 \(BDE\) 内，所以 \(GE\subset\) 平面 \(BDE\).平面 \(BDE\) 与平面 \(ABCD\) 的交线为 \(BD\)，而 \(A\notin BD\)，故 \(A\notin\) 平面 \(BDE\)，从而 \(AF\not\subset\) 平面 \(BDE\).因此 \(AF\parallel\) 平面 \(BDE\).

\item 又因为 \(EF\parallel CG\)，且 \(EF=CG=1\)，所以四边形 \(CEFG\) 是平行四边形.结合 \(CE=EF=1\)，它是菱形，因此其对角线 \(CF\perp EG\).正方形中 \(BD\perp AC\)，且正方形 \(ABCD\) 所在平面与四边形 \(ACEF\) 所在平面垂直，所以 \(BD\perp\) 平面 \(ACEF\)，从而 \(BD\perp CF\).直线 \(BD\) 与 \(EG\) 在平面 \(BDE\) 内相交，且 \(CF\) 同时垂直于它们，所以 \(CF\perp\) 平面 \(BDE\).

\end{enumerate}
\end{solution}

\begin{problem}
设函数 \( f(x) = \frac{a}{3} x^3 + bx^2 + cx + d\) （\(a > 0\)），且方程 \( f'(x) - 9x = 0 \) 的两个根分别为 1，4.

\begin{enumerate}
\item 当 \(a = 3\) 且曲线 \(y = f(x)\) 过原点时，求 \(f(x)\) 的解析式；
\item 若 \( f(x) \) 在 \( (-\infty,+\infty) \) 内无极值点，求 a 的取值范围.
\end{enumerate}
\end{problem}

\begin{solution}
由
\[
f'(x)=ax^2+2bx+c
\]
方程 \(f'(x)-9x=0\) 为
\[
ax^2+(2b-9)x+c=0
\]
它的两个根为 \(1,4\)，由韦达定理得
\(\frac{9-2b}{a}=5\)，\(\frac{c}{a}=4\)
即 \(2b=9-5a\)，\(c=4a\).

\begin{enumerate}
\item 当 \(a=3\) 时，\(b=-3\)，\(c=12\).又曲线过原点，故 \(d=f(0)=0\)，所以
\[
f(x)=x^3-3x^2+12x
\]

\item 由 \(2b=9-5a\)，\(c=4a\)，得
\[
f'(x)=ax^2+(9-5a)x+4a
\]
因 \(a>0\)，要使 \(f(x)\) 在整个实数范围内无极值点，需且只需 \(f'(x)\geqslant0\) 对一切实数 \(x\) 成立，即该二次函数的判别式不大于 \(0\).因此
\[
\Delta=(9-5a)^2-16a^2=9(a-1)(a-9)\leqslant0
\]
结合 \(a>0\)，得 \(1\leqslant a\leqslant9\)，所以 \(a\in[1,9]\).
\end{enumerate}
\end{solution}

\begin{problem}
已知椭圆 \(C\) 的左，右焦点坐标分别是 \((- \sqrt{2}, 0)\)，\((\sqrt{2}, 0)\)，离心率是 \(\frac{\sqrt{6}}{3}\). 直线 \(y = t\) 与椭圆 \(C\) 交于不同的两点 \(M\)，\(N\)，以线段 \(MN\) 为直径作圆 \(P\)，圆心为 \(P\).

\begin{enumerate}
\item 求椭圆 \(C\) 的方程；
\item 若圆 \(P\) 与 \(x\) 轴相切，求圆心 \(P\) 的坐标；
\item 设 \(Q(x,y)\) 是圆 \(P\) 上的动点，当 \(t\) 变化时，求 \(y\) 的最大值.
\end{enumerate}
\end{problem}

\begin{solution}
\begin{enumerate}
\item 设椭圆的半长轴、半短轴和焦半距分别为 \(a\)、\(b\)、\(c\).由焦点坐标知 \(c=\sqrt{2}\)，又离心率为 \(\frac{c}{a}=\frac{\sqrt{6}}{3}\)，所以 \(a=\sqrt{3}\).于是 \(b^2=a^2-c^2=3-2=1\)，椭圆方程为
\[
\frac{x^2}{3}+y^2=1
\]

\item 直线 \(y=t\) 与椭圆交于不同两点，故 \(-1<t<1\)，且
\(M\left(\sqrt{3(1-t^2)},t\right)\)，\(N\left(-\sqrt{3(1-t^2)},t\right)\).
所以圆心 \(P=(0,t)\)，半径为 \(r=\sqrt{3(1-t^2)}\).圆与 \(x\) 轴相切等价于 \(|t|=r\)，即
\[
t^2=3(1-t^2)
\]
解得 \(t=\pm\frac{\sqrt{3}}{2}\).因此圆心坐标为 \(\left(0,\pm\frac{\sqrt{3}}{2}\right)\).

\item 圆 \(P\) 的方程为
\[
x^2+(y-t)^2=3(1-t^2)
\]
对固定的 \(t\)，圆上点的纵坐标满足
\[
y\leqslant t+\sqrt{3(1-t^2)}
\]
令 \(u=t\)，\(v=\sqrt{1-t^2}\)，则 \(u^2+v^2=1\)，且 \(v>0\).由柯西不等式，
\[
y\leqslant u+\sqrt{3}v\leqslant\sqrt{1^2+(\sqrt{3})^2}\sqrt{u^2+v^2}=2
\]
当 \(u=\frac{1}{2}\)，\(v=\frac{\sqrt{3}}{2}\)，即 \(t=\frac{1}{2}\) 且 \(x=0\) 时取得等号.因此 \(y\) 的最大值为 \(2\).
\end{enumerate}
\end{solution}

\begin{problem}
已知集合 \( S_{n} = \{X \mid X = (x_{1}, x_{2}, \ldots, x_{n}), x_{i} \in \{0, 1\}, i = 1, 2, \cdots, n\} \)（\( n \geqslant 2 \)），对于 \( A = (a_{1}, a_{2}, \cdots, a_{n}) \)，\( B = (b_{1}, b_{2}, \cdots, b_{n}) \in S_{n} \)，定义 \(A\) 与 \(B\) 的差为 \( A - B = (|a_{1} - b_{1}|, |a_{2} - b_{2}|, \cdots, |a_{n} - b_{n}|) \)； \(A\) 与 \(B\) 之间的距离为 \( d(A, B) = \sum_{i=1}^{n} |a_{i} - b_{i}| \).

\begin{enumerate}
\item 当 \(n = 5\) 时，设 \(A = (0, 1, 0, 0, 1)\)，\(B = (1, 1, 1, 0, 0)\)，求 \(A - B\)，\(d(A, B)\)；
\item 证明： \(\forall A\)，\(B\)，\(C \in S_n\)，有 \(A - B \in S_n\)，且 \(d(A - C, B - C) = d(A, B)\)；
\item 证明： \(\forall A\)，\(B\)，\(C \in S_n\)，\(d(A, B)\)，\(d(A, C)\)，\(d(B, C)\) 三个数中至少有一个是偶数.
\end{enumerate}
\end{problem}

\begin{solution}
\begin{enumerate}
\item
\[
A-B=(|0-1|,|1-1|,|0-1|,|0-0|,|1-0|)=(1,0,1,0,1)
\]
所以
\[
d(A,B)=1+0+1+0+1=3
\]

\item 因为 \(a_i\)、\(b_i\in\{0,1\}\)，所以 \(|a_i-b_i|\in\{0,1\}\)，从而
\[
A-B=(|a_1-b_1|,\ldots,|a_n-b_n|)\in S_n
\]
对每个 \(i\)，若 \(c_i=0\)，则
\[
\bigl||a_i-c_i|-|b_i-c_i|\bigr|=|a_i-b_i|
\]
若 \(c_i=1\)，则 \(|a_i-c_i|=1-a_i\)，\(|b_i-c_i|=1-b_i\)，所以
\[
\bigl||a_i-c_i|-|b_i-c_i|\bigr|=|(1-a_i)-(1-b_i)|=|a_i-b_i|
\]
对 \(i=1\)，\(2\)，\(\ldots\)，\(n\) 求和，得到
\[
d(A-C,B-C)=\sum_{i=1}^{n}|a_i-b_i|=d(A,B)
\]

\item 对任意一个坐标 \(i\)，若 \(a_i\)、\(b_i\)、\(c_i\) 全相等，则三项 \(|a_i-b_i|\)、\(|a_i-c_i|\)、\(|b_i-c_i|\) 全为 \(0\)；若不全相等，由于它们只能取 \(0,1\)，则三者中恰有一个与另外两个不同，于是上述三项中恰有两个为 \(1\).因此每个坐标对
\[
|a_i-b_i|+|a_i-c_i|+|b_i-c_i|
\]
的贡献都是偶数，从而
\[
d(A,B)+d(A,C)+d(B,C)
\]
是偶数.若三个距离都为奇数，其和将为奇数，矛盾.因此三个数中至少有一个是偶数.
\end{enumerate}
\end{solution}
